\SourceLeaf{43}\begin{center}
{\Large DRITTER ABSCHNITT}\par\bigskip
{\small ÜBER}\par\bigskip
{\Large POTENZRESTE.}
\end{center}

\begin{center}\rule{0.09\linewidth}{0.3pt}\end{center}

\begin{center}\small\emph{Die Reste der Glieder einer mit der Eins beginnenden geometrischen Folge bilden eine periodische Folge.}\end{center}

\Article{45}

\textsc{Satz.} \emph{In jeder geometrischen Folge $1,a,aa,a^3$ usw. gibt es außer dem ersten Glied $1$ noch ein Glied $a^t$, das nach einem Modul $p$, der zu $a$ teilerfremd ist, der Eins kongruent ist und den Exponenten $t<p$ hat.}

\emph{Beweis.} Da der Modul $p$ zu $a$ und damit zu jeder Potenz von $a$ teilerfremd ist, ist kein Glied der Folge $\equiv0\ (\mathrm{mod.}p)$. Jedes ist vielmehr zu einer der Zahlen $1,2,3....p-1$ kongruent. Da ihre Anzahl $p-1$ ist, können offenbar nicht alle Glieder verschiedene kleinste Reste haben, wenn man mehr als $p-1$ Glieder der Folge betrachtet. Unter den Gliedern $1,a,aa,a^3....a^{p-1}$ sind also mindestens zwei kongruent. Sei $a^m\equiv a^n$ und $m>n$. Durch Division mit $a^n$ ergibt sich $a^{m-n}\equiv1$ (Artikel 22), wobei $m-n<p$ und $>0$ ist.\par
\noindent\mbox{\emph{Q. E. D.}}

\emph{Beispiel.} In der Folge $2,4,8$ usw. ist das erste Glied, das nach dem Modul $13$ der Eins kongruent ist, $2^{12}=4096$. Nach dem Modul $23$ ergibt sich dagegen in derselben Folge $2^{11}=2048\equiv1$. Ebenso ist die sechste Potenz von $5$, nämlich $15625$, nach dem Modul $7$ der Eins kongruent, die fünfte Potenz, $3125$, dagegen nach $11$. In manchen Fällen ist also schon eine Potenz mit kleinerem Exponenten als $p-1$ der Eins kongruent; in anderen muss man bis zur Potenz mit dem Exponenten $p-1$ gehen.

\SourceLeaf{44}\Article{46}

Setzt man die Folge über ein der Eins kongruentes Glied hinaus fort, so treten dieselben Reste wie am Anfang erneut auf. Ist nämlich $a^t\equiv1$, so ist $a^{t+1}\equiv a,a^{t+2}\equiv aa$ usw., bis man zum Glied $a^{2t}$ kommt, dessen kleinster Rest wieder $\equiv1$ ist und die \emph{Periode} der Reste neu beginnt. Man hat also eine Periode mit $t$ Resten, die nach ihrem Ende stets von vorn wiederholt wird. Andere Reste als die in dieser Periode enthaltenen können in der ganzen Folge nicht vorkommen. Allgemein ist $a^{mt}\equiv1$ und $a^{mt+n}\equiv a^n$, was sich in unserer Schreibweise so ausdrückt:

\[\text{\emph{Ist }}r\equiv\rho\ (mod.\,t),\quad\text{\emph{so ist}}\quad a^r\equiv a^\rho\ (mod.\,p).\]

\Article{47}

Dieser Satz ermöglicht es, Reste von Potenzen mit beliebig großen Exponenten rasch zu bestimmen, sobald eine der Eins kongruente Potenz bekannt ist. Sucht man \emph{zum Beispiel} den Rest bei Division von $3^{1000}$ durch $13$, so ist wegen $3^3\equiv1\ (\mathrm{mod.}\ 13)$ auch $t\equiv3$. Da $1000\equiv1\ (\mathrm{mod.}\ 3)$ ist, folgt $3^{1000}\equiv3\ (\mathrm{mod.}\ 13)$.

\Article{48}

Ist $a^t$ die \emph{niedrigste} der Eins kongruente Potenz, abgesehen von $a^0=1$, was wir hier außer Betracht lassen, so sind jene $t$ Glieder, die die Periode der Reste bilden, alle verschieden, wie man aus dem Beweis von Artikel 45 leicht erkennt. Dann lässt sich der Satz von Artikel 46 umkehren: Ist $a^m\equiv a^n\ (\mathrm{mod.}p)$, so ist $m\equiv n\ (\mathrm{mod.}t)$. Wären nämlich $m,n$ nach dem Modul $t$ inkongruent, so wären ihre kleinsten Reste $\mu,\nu$ verschieden. Es ist aber $a^\mu\equiv a^m,a^\nu\equiv a^n$, also $a^\mu\equiv a^\nu$. \emph{Das heißt}, nicht alle Potenzen unter $a^t$ wären inkongruent, entgegen der Voraussetzung.

Ist also $a^k\equiv1\ (\mathrm{mod.}p)$, so ist $k\equiv0\ (\mathrm{mod.}t)$, \emph{das heißt}, $k$ ist durch $t$ teilbar.

Bisher haben wir beliebige Moduln betrachtet, sofern sie zu $a$ teilerfremd sind. Nun betrachten wir Primzahlmoduln gesondert und bauen später auf dieser Grundlage eine allgemeinere Untersuchung auf.

\SourceLeaf{45}\begin{center}\small\emph{Zunächst werden Moduln betrachtet, die Primzahlen sind.}\end{center}

\Article{49}

\textsc{Satz.} \emph{Ist $p$ eine Primzahl, die $a$ nicht teilt, und ist $a^t$ die niedrigste Potenz von $a$, die nach dem Modul $p$ der Eins kongruent ist, so ist der Exponent $t$ entweder $=p-1$ oder ein Teiler dieser Zahl.}

Man vergleiche die Beispiele in Artikel 45.

\emph{Beweis.} Es ist bereits gezeigt, dass $t$ entweder $=p-1$ oder $<p-1$ ist. Es bleibt zu beweisen, dass im letzteren Fall $t$ stets ein Teiler von $p-1$ ist.

I. Man nehme die kleinsten positiven Reste aller Glieder $1,a,aa....a^{t-1}$ und bezeichne sie mit $\alpha,\alpha',\alpha''$ usw., sodass $\alpha=1,\alpha'\equiv a,\alpha''\equiv aa$ usw. gilt. Sie sind offenbar alle verschieden. Ergäben nämlich zwei Glieder $a^m,a^n$ denselben Rest, so wäre, wenn man $m>n$ annimmt, $a^{m-n}\equiv1$ und $m-n<t$, \mbox{\emph{Q. E. A.}}, da keine niedrigere Potenz als $a^t$ der Eins kongruent ist (\emph{nach Voraussetzung}). Ferner gehören alle Zahlen $\alpha,\alpha',\alpha''$ usw. zur Reihe $1,2,3...p-1$, schöpfen sie aber nicht aus, da $t<p-1$ ist. Die Gesamtheit aller $\alpha,\alpha',\alpha''$ usw. bezeichnen wir mit $(A)$. Somit enthält $(A)$ genau $t$ Glieder.

II. Man wähle eine beliebige Zahl $\beta$ aus $1,2,3...p-1$, die in $(A)$ fehlt, multipliziere $\beta$ mit allen $\alpha,\alpha',\alpha''$ usw. und nenne die daraus entstehenden kleinsten Reste $\beta,\beta',\beta''$ usw. Ihre Anzahl ist ebenfalls $t$. Diese Reste sind sowohl untereinander als auch von allen $\alpha,\alpha',\alpha''$ usw. verschieden. Wäre die \emph{erste} Behauptung falsch, so wäre $\beta a^m\equiv\beta a^n$ und nach Division durch $\beta$ auch $a^m\equiv a^n$, entgegen dem eben Bewiesenen. Wäre die \emph{zweite} falsch, so wäre $\beta a^m\equiv a^n$. Für $m<n$ folgte $\beta\equiv a^{n-m}$, \emph{das heißt}, $\beta$ wäre zu einer der Zahlen $\alpha,\alpha',\alpha''$ usw. kongruent, entgegen der Voraussetzung. Für $m>n$ folgt durch Multiplikation mit $a^{t-m}$ die Kongruenz $\beta a^t\equiv a^{t+n-m}$ oder wegen $a^t\equiv1$ auch $\beta\equiv a^{t+n-m}$, derselbe Widerspruch. Die Gesamtheit aller $\beta,\beta',\beta''$ usw., deren Anzahl $=t$ ist, bezeichnen wir mit $(B)$. Damit haben wir bereits $2t$ Zahlen aus $1,2,3...p-1$. Umfassen also $(A)$ und $(B)$ alle diese Zahlen, so ist $\frac{p-1}{2}=t$ und der Satz bewiesen.

III. Fehlen noch Zahlen, so sei $\gamma$ eine davon. Man multipliziere damit alle $\alpha,\alpha',\alpha''$ usw.; die kleinsten Reste der Produkte seien $\gamma,\gamma',\gamma''$ usw. Ihre Gesamtheit werde mit $(C$ bezeichnet. Dann enthält $(C)$ genau $t$ Zahlen aus $1,2,3...p-1$, die sowohl untereinander als auch von den in $(A)$ und $(B)$
\SourceLeaf{46}
enthaltenen Zahlen verschieden sind. Die ersten beiden Behauptungen werden wie in II bewiesen, die dritte folgendermaßen: Wäre $\gamma a^m\equiv\beta a^n$, so folgte $\gamma\equiv\beta a^{n-m}$ oder $\equiv\beta a^{t+n-m}$, je nachdem, ob $m<n$ oder $>n$ ist. In beiden Fällen wäre $\gamma$ zu einer Zahl aus $(B)$ kongruent, entgegen der Voraussetzung. Man hat also $3t$ Zahlen aus $1,2,3...p-1$. Fehlen keine mehr, so ist $t=\frac{p-1}{3}$ und der Satz bewiesen.

IV. Fehlen immer noch Zahlen, so geht man auf dieselbe Weise zu einer vierten Gesamtheit $(D)$ über usw. Da die Anzahl der Zahlen $1,2,3....p-1$ endlich ist, müssen sie schließlich alle erfasst sein. Ihre Anzahl ist daher ein Vielfaches von $t$, also ist $t$ ein Teiler von $p-1$. \mbox{\emph{Q. E. D.}}

\Needspace{6\baselineskip}
\begin{center}\small\emph{Fermats Satz.}\end{center}

\Article{50}

Da also $\frac{p-1}{t}$ eine ganze Zahl ist, erhält man durch Erheben beider Seiten der Kongruenz $a^t\equiv1$ zur Potenz mit dem Exponenten $\frac{p-1}{t}$ die Kongruenz $a^{p-1}\equiv1$. \emph{Das heißt, $a^{p-1}-1$ ist stets durch $p$ teilbar, wenn $p$ eine Primzahl ist, die $a$ nicht teilt.}

Dieser Satz verdient wegen seiner Eleganz und seiner außerordentlichen Nützlichkeit volle Aufmerksamkeit. Nach seinem Entdecker wird er gewöhnlich \emph{Fermats Satz} genannt. Siehe \emph{Fermatii Opera Mathem. Tolosae} 1679, \emph{fol. S.} 163. Der Entdecker fügte keinen Beweis hinzu, erklärte aber, einen zu besitzen. Der berühmte Euler veröffentlichte als Erster einen Beweis in der Abhandlung \emph{Theorematum quorundam ad numeros primos spectantium demonstratio, Comm. Acad. Petrop. T.} VIII\begingroup\renewcommand{\thefootnote}{*)}\footnote{In einer früheren Abhandlung hatte dieser große Mann sein Ziel noch nicht erreicht: \emph{Comm. Petr. T.} VI, \emph{S.} 106. --- In dem berühmten Streit zwischen Maupertuis und Konig, der vom Prinzip der kleinsten Wirkung ausging, sich aber bald auf andere Gegenstände ausdehnte, erklärte K\"onig, eine eigenhändige Niederschrift von Leibniz zu besitzen, die einen mit Eulers Beweis völlig übereinstimmenden Beweis dieses Satzes enthalte: \emph{Appel au public. S.} $106$. Wir wollen diesem Zeugnis zwar nicht die Glaubwürdigkeit absprechen; gewiss ist aber, dass Leibniz seine Entdeckung nie veröffentlichte. Vergleiche \emph{Hist. de l'Ac. de Prusse, A.} 1750, \emph{S.} 530.}\endgroup. Dieser Beweis beruht auf der Entwicklung der Potenz $(a+1)^p$. Aus der Form der Koeffizienten folgt leicht, dass $(a+1)^p-a^p-1$ stets durch $p$ teilbar ist. Daher ist $(a+1)^p-(a+1)$ durch $p$ teilbar, wenn $a^p-a$ durch $p$ teilbar ist. Da nun $1^p-1$ stets durch $p$ teilbar ist, gilt dies auch für $2^p-2$, damit auch für $3^p-3$ usw. und allgemein für $a^p-a$. Teilt also $p$ die Zahl $a$ nicht, so ist auch $a^{p-1}-1$ durch $p$ teilbar. Dies genügt, um die Art der Methode zu zeigen. Der angesehene Lambert gab einen ähnlichen Beweis in \emph{Actis Erudit.} 1769,
\SourceLeaf{47}
\emph{S.} 109. Da aber die Entwicklung einer Binomialpotenz der Zahlentheorie etwas fremd zu sein schien, suchte der berühmte Euler einen anderen Beweis. Er findet sich in \emph{Comment. nov. Petr. T.} VII, \emph{S.} 70, und stimmt völlig mit dem Beweis überein, den wir im vorigen Artikel gegeben haben. Später werden sich uns noch weitere Beweise bieten. Hier dürfen wir noch einen hinzufügen, der auf ähnlichen Grundsätzen wie Eulers erster Beweis beruht. Der folgende Satz, von dem unser Satz nur ein Sonderfall ist, wird auch bei anderen Untersuchungen weiter unten verwendet werden.

\Article{51}

\emph{Das Polynom $a+b+c+\mathit{usw.}$ ist, zur Potenz mit dem Exponenten $p$ erhoben, nach dem Modul $p$}

\[\equiv a^p+b^p+c^p+\mathit{usw.}\]

\noindent\emph{sofern $p$ eine Primzahl ist.}

\emph{Beweis.} Die Potenz mit dem Exponenten $p$ des Polynoms $a+b+c+\text{usw.}$ besteht aus Gliedern der Form $\varkappa a^\alpha b^\beta c^\gamma$ usw., wobei $\alpha+\beta+\gamma\text{ usw.}=p$ ist. Dabei bezeichnet $\varkappa$ die Anzahl der Permutationen von $p$ Gegenständen, von denen $\alpha,\beta,\gamma$ usw. jeweils $=a,b,c$ usw. sind. Nach Artikel 41 ist diese Anzahl stets durch $p$ teilbar, außer wenn alle Gegenstände gleich sind, \emph{das heißt}, wenn eine der Zahlen $\alpha,\beta,\gamma$ usw. gleich $=p$ und die übrigen gleich $=0$ sind. Somit sind alle Glieder von $(a+b+c+\text{usw.})^p$ außer $a^p,b^p,c^p$ usw. durch $p$ teilbar. Bei einer Kongruenz nach dem Modul $p$ dürfen sie also wegfallen, und man erhält

\[(a+b+c+\mathrm{usw.})^p\equiv a^p+b^p+c^p+\mathrm{usw.}\quad\text{\emph{Q. E. D.}}\]

Setzt man nun alle Größen $a,b,c$ usw. $=1$ und ihre Anzahl $=k$, so ergibt sich $k^p\equiv k$, wie im vorigen Artikel.

\Needspace{6\baselineskip}
\begin{center}\emph{Wie viele Zahlen zu Perioden gehören, deren Gliederzahl ein gegebener Teiler von $p-1$ ist.}\end{center}

\Article{52}

Da nur Teiler von $p-1$ Exponenten der niedrigsten Potenzen sein können, zu denen Zahlen erhoben werden, um der Eins kongruent zu werden, stellt sich die Frage, ob alle Teiler von $p-1$ dafür geeignet sind. Ferner ist zu fragen, wie viele Zahlen zu jedem Exponenten gehören, wenn man alle nicht durch $p$ teilbaren Zahlen nach dem Exponenten ihrer niedrigsten der Eins kongruenten Potenz einteilt. Zunächst genügt es offenbar, die positiven Zahlen
\SourceLeaf{48}
von $1$ bis $p-1$ zu betrachten. Kongruente Zahlen müssen nämlich zur gleichen niedrigsten Potenz erhoben werden, um der Eins kongruent zu werden. Daher gehört jede Zahl zu demselben Exponenten wie ihr kleinster positiver Rest. Wir müssen also ermitteln, wie die Zahlen $1,2,3\GaussFourDots p-1$ in dieser Hinsicht auf die einzelnen Faktoren von $p-1$ verteilt sind. Ist $d$ einer der Teiler von $p-1$, zu denen auch $1$ und $p-1$ gehören, so bezeichnen wir der Kürze halber mit $\psi d$ die Anzahl der positiven Zahlen unter $p$, deren niedrigste der Eins kongruente Potenz den Exponenten $d$ hat.

\Article{53}

Um die Untersuchung leichter verständlich zu machen, geben wir ein Beispiel. Für $p=19$ verteilen sich die Zahlen $1,2,3\GaussFourDots18$ auf die Teiler von $18$ folgendermaßen:

\[
\begin{array}{r|l}
1 & \phantom{1}1.\\
2 & 18.\\
3 & \phantom{1}7,\ 11.\\
6 & \phantom{1}8,\ 12.\\
9 & \phantom{1}4,\ 5.\ 6,\ 9,\ 16,\ 17.\\
18 & \phantom{1}2,\ 3,\ 10,\ 13.\ 14,\ 15.
\end{array}
\]

In diesem Fall ist also $\psi1=1$, $\psi2=1$, $\psi3=2$, $\psi6=2$, $\psi9=6$, $\psi18=6$. Bei einiger Aufmerksamkeit erkennt man, dass zu jedem Exponenten ebenso viele Zahlen gehören, wie es nicht größere und zu ihm teilerfremde Zahlen gibt. Mit der Schreibweise von Artikel 39 ist also zumindest in diesem Fall $\psi d=\phi d$. Dass dies allgemein gilt, beweisen wir so.

I. Es sei eine Zahl $a$ gegeben, die zum Exponenten $d$ \emph{gehört}, \emph{das heißt}, deren Potenz mit dem Exponenten $d$ der Eins kongruent ist, während alle niedrigeren Potenzen dazu inkongruent sind. Dann haben auch alle ihre Potenzen $aa,a^3,a^4\GaussFourDots a^d$ oder deren kleinste Reste die erste Eigenschaft: Ihre Potenzen mit dem Exponenten $d$ sind der Eins kongruent. Anders ausgedrückt sind die kleinsten Reste der Zahlen $a,aa,a^3\GaussFourDots a^d$, die alle verschieden sind, Wurzeln der Kongruenz $x^d\equiv1$. Da diese nicht mehr als $d$ verschiedene Wurzeln haben kann, gibt es außer den kleinsten Resten der Zahlen $a,aa,a^3\GaussFourDots a^d$ keine weiteren Zahlen zwischen $1$ und $p-1$ einschließlich, deren Potenzen mit dem Ex\SourceLeaf{49}ponenten $d$ der Eins kongruent sind. Alle zum Exponenten $d$ gehörenden Zahlen finden sich also unter den kleinsten Resten der Zahlen $a,aa,a^3\GaussFourDots a^d$. Welche es sind und wie viele, bestimmt man folgendermaßen. Ist $k$ zu $d$ teilerfremd, so ist keine Potenz von $a^k$ mit einem Exponenten $<d$ der Eins kongruent. Sei nämlich $\frac{1}{k}\,(\operatorname{mod.}d)\equiv m$ (siehe Artikel 31); dann ist $a^{km}\equiv a$. Wäre nun die Potenz mit dem Exponenten $e$ von $a^k$ der Eins kongruent und $e<d$, so wäre auch $a^{kme}\equiv1$ und damit $a^e\equiv1$, entgegen der Voraussetzung. Also gehört der kleinste Rest von $a^k$ offenbar zum Exponenten $d$. Hat dagegen $k$ einen gemeinsamen Teiler $\delta$ mit $d$, so gehört der kleinste Rest von $a^k$ nicht zum Exponenten $d$. Denn schon die Potenz mit dem Exponenten ${\frac{d}{\delta}}$ ist der Eins kongruent, da $\frac{kd}{\delta}$ durch $d$ teilbar ist, also $\equiv0\,(\operatorname{mod.}d)$ und damit $a^{\frac{kd}{\delta}}\equiv1$ gilt. Folglich gehören zum Exponenten $d$ ebenso viele Zahlen, wie unter $1,2,3\GaussFourDots d$ zu $d$ teilerfremde Zahlen vorkommen. Dabei ist aber zu beachten, dass wir vorausgesetzt haben, bereits eine Zahl $a$ zu besitzen, die zum Exponenten $d$ gehört. Es bleibt also fraglich, ob zu einem Exponenten vielleicht überhaupt keine Zahl gehört. Der Schluss beschränkt sich daher darauf, dass $\psi d$ entweder $=0$ oder $=\phi d$ ist.

\Article{54}

II. Alle Teiler von $p-1$ seien nun $d,d',d''$ usw. Da alle Zahlen $1,2,3\GaussFourDots p-1$ auf diese verteilt sind, ist

\[\psi d+\psi d'+\psi d''+\text{usw.}=p-1\]

In Artikel 40 haben wir aber bewiesen, dass

\[\phi d+\phi d'+\phi d''+\text{usw.}=p-1\]

Nach dem vorigen Artikel ist $\psi d$ entweder gleich oder kleiner als $\phi d$, kann aber nicht größer sein; ebenso gilt dies für $\psi d'$ und $\phi d'$ usw. Wäre nun eines der Glieder $\psi d,\psi d',\psi d''$ usw. kleiner als das entsprechende Glied unter $\phi d,\phi d',\phi d''$, oder wären es mehrere, so könnten die beiden Summen nicht gleich sein. Wir schließen also schließlich, dass $\psi d$ \emph{stets gleich $\phi d$ ist} und daher nicht von der Größe von $p-1$ abhängt.

\Article{55}

Besondere Aufmerksamkeit verdient der Sonderfall des vorigen\SourceLeaf{50} Satzes, dass es nämlich \emph{stets Zahlen gibt, deren Potenzen mit kleinerem Exponenten als $p-1$ der Eins nicht kongruent sind}. Und zwar gibt es zwischen $1$ und $p-1$ ebenso viele davon, wie es unter $p-1$ zu $p-1$ teilerfremde Zahlen gibt. Der Beweis dieses Satzes ist keineswegs so selbstverständlich, wie es auf den ersten Blick scheinen mag. Wegen seiner Bedeutung dürfen wir noch einen etwas anderen Beweis hinzufügen, da verschiedene Methoden oft viel dazu beitragen, schwierigere Dinge verständlich zu machen. Man zerlege $p-1$ in seine Primfaktoren, sodass $p-1=a^\alpha b^\beta c^\gamma$ usw. ist, wobei $a,b,c$ usw. verschiedene Primzahlen bezeichnen. Zum Beweis zeigen wir nun Folgendes:

I. Es lässt sich stets eine Zahl $A$, oder auch mehrere, finden, die zum Exponenten $a^\alpha$ gehört; ebenso Zahlen $B,C$ usw., die zu den Exponenten $b^\beta,c^\gamma$ usw. jeweils gehören.

II. Das Produkt aller Zahlen $A,B,C$ usw. oder dessen kleinster Rest gehört zum Exponenten $p-1$. Dies beweisen wir folgendermaßen.

I. Sei $g$ eine Zahl aus $1,2,3\ldots p-1$, welche die Kongruenz $x^{\frac{p-1}{a}}\equiv1\,(\operatorname{mod.}p)$ \emph{nicht} erfüllt. Denn nicht alle diese Zahlen können eine Kongruenz erfüllen, deren Grad $<p-1$ ist. Ich behaupte nun: Ist die Potenz mit dem Exponenten ${\frac{p-1}{a^\alpha}}$ von $g$ $\equiv h$, so gehört diese Zahl oder ihr kleinster Rest zum Exponenten $a^\alpha$.

Denn die Potenz mit dem Exponenten $a^{\alpha}$ von $h$ ist offenbar zur Potenz mit dem Exponenten $p-1$ von $g$, \emph{also} zur Eins kongruent. Die Potenz mit dem Exponenten $a^{\alpha-1}$ von $h$ ist dagegen zur Potenz mit dem Exponenten ${\frac{p-1}{a}}$ von $g$ kongruent, \emph{also} zur Eins inkongruent. Noch weniger können die Potenzen mit den Exponenten $a^{\alpha-2},\ a^{\alpha-3}$ usw. von $h$ der Eins kongruent sein. Der Exponent der niedrigsten der Eins kongruenten Potenz von $h$, also der Exponent, zu dem $h$ gehört, muss aber ein Teiler von $a^\alpha$ sein (Artikel 48). Da $a^\alpha$ nur durch sich selbst und niedrigere Potenzen von $a$ teilbar ist, muss $a^\alpha$ der Exponent sein, zu dem $h$ gehört. \mbox{\emph{Q. E. D.}} Ebenso zeigt man, dass es Zahlen gibt, die zu den Exponenten $b^\beta,c^\gamma$ usw. gehören.

II. Nehmen wir an, das Produkt aller $A,B,C$ usw. gehöre nicht zum Exponenten $p-1$, sondern zu einem kleineren $t$. Dann teilt $t$ die Zahl $p-1$ (Artikel 48), also ist $\frac{p-1}{t}$ eine ganze Zahl größer als eins. Offenbar ist dieser Quotient entweder eine der Primzahlen $a,b,c$ usw. oder zumindest durch eine davon teilbar (Artikel 17), \emph{zum Beispiel} durch $a$. Für die übrigen verläuft der Schluss ebenso.

\SourceLeaf{51}
Somit teilt $t$ die Zahl $\frac{p-1}{a}$, und das Produkt $ABC$ usw. ist auch dann der Eins kongruent, wenn es zur Potenz mit dem Exponenten ${\frac{p-1}{a}}$ erhoben wird (Artikel 46). Aber auch jede einzelne der Zahlen $B,C$ usw., ausgenommen $A$, ist zur Potenz mit dem Exponenten ${\frac{p-1}{a}}$ erhoben der Eins kongruent. Denn die Exponenten $b^\beta,c^\gamma$ usw., zu denen sie gehören, teilen $\frac{p-1}{a}$. Daher ist
\[A^{\frac{p-1}{a}}B^{\frac{p-1}{a}}C^{\frac{p-1}{a}}\text{ usw.}\equiv A^{\frac{p-1}{a}}\equiv1.\]
Der Exponent, zu dem $A$ gehört, muss also $\frac{p-1}{a}$ teilen (Artikel 48), \emph{das heißt}, $\frac{p-1}{a^{\alpha+1}}$ muss ganzzahlig sein. Aber $\frac{p-1}{a^{\alpha+1}}=\frac{b^\beta c^\gamma\text{ usw.}}{a}$ kann keine ganze Zahl sein (Artikel 15). Unsere Annahme kann daher nicht bestehen; \emph{das heißt}, das Produkt $ABC$ usw. gehört tatsächlich zum Exponenten $p-1$. \mbox{\emph{Q. E. D.}}

Der zweite Beweis scheint etwas länger als der erste, der erste dagegen weniger direkt als der zweite zu sein.

\Article{56}
Dieser Satz zeigt besonders deutlich, wie viel Vorsicht in der Zahlentheorie oft nötig ist, damit wir nicht als sicher annehmen, was nicht sicher ist. Der angesehene Lambert erwähnt diesen Satz in der bereits angeführten Abhandlung \emph{Acta Erudit.} 1769, \emph{S.} 127, geht aber nicht einmal darauf ein, dass ein Beweis nötig ist. Niemand außer dem großen Euler hat einen Beweis versucht: \emph{Comment. nov. Ac. Petrop. T.} XVIII zum Jahr 1773, \emph{Demonstrationes circa residua ex divisione potestatum per numeros primos resultantia}, S. 85 ff. Man vergleiche besonders Artikel 37, in dem er die Notwendigkeit eines Beweises ausführlicher bespricht. Der Beweis dieses scharfsinnigen Mannes hat jedoch zwei Lücken. Erstens setzt er in Artikel 31 ff. stillschweigend voraus, dass die Kongruenz $x^n\equiv1$, wenn man seine Überlegungen in unsere Zeichen überträgt, tatsächlich $n$ verschiedene Wurzeln hat, obwohl vorher nur gezeigt war, dass sie nicht \emph{mehr} haben kann. Zweitens hat er die Formel in Artikel 34 nur durch Induktion hergeleitet.

\Needspace{6\baselineskip}
\begin{center}\emph{Primitive Wurzeln, Basen, Indizes.}\end{center}

\Article{57}
Mit dem berühmten Euler nennen wir die zum Exponenten $p-1$ gehörenden Zahlen \emph{primitive Wurzeln}. Ist also $a$ eine primitive Wurzel, so sind die kleinsten Reste der Potenzen $a,aa,a^3\ldots a^{p-1}$
\SourceLeaf{52}
alle verschieden. Sie müssen daher sämtliche Zahlen $1,2,3,\ldots p-1$ enthalten, da diese ebenso zahlreich sind wie jene kleinsten Reste. Das heißt, jede nicht durch $p$ teilbare Zahl ist zu einer Potenz von $a$ kongruent. Diese bemerkenswerte Eigenschaft ist sehr nützlich und kann die Rechenoperationen mit Kongruenzen erheblich erleichtern, ähnlich wie die Einführung der Logarithmen die gewöhnliche Arithmetik erleichtert. Wir wählen eine beliebige primitive Wurzel $a$ als \emph{Basis}, auf die wir alle nicht durch $p$ teilbaren Zahlen beziehen. Ist $a^e\equiv b\,(\operatorname{mod.}p)$, so nennen wir $e$ den \emph{Index} von $b$. \emph{Zum Beispiel} entsprechen für den Modul $19$, wenn die primitive Wurzel $2$ als Basis gewählt wird, einander:

\begin{center}
\begingroup\small\setlength{\tabcolsep}{1.5pt}
\begin{tabular}{@{}l@{\quad}*{18}{r}@{}}
Zahlen & 1. & 2. & 3. & 4. & 5. & 6. & 7. & 8. & 9. & 10. & 11. & 12. & 13. & 14. & 15. & 16. & 17. & 18.\\
Indizes & 0. & 1. & 13. & 2. & 16. & 14. & 6. & 3. & 8. & 17. & 12. & 15. & 5. & 7. & 11. & 4. & 10. & 9.
\end{tabular}
\endgroup
\end{center}

Bei gleicher Basis gehören zu jeder Zahl offenbar mehrere Indizes, die aber alle nach dem Modul $p-1$ kongruent sind. Sprechen wir von Indizes, so gelten daher diejenigen als gleichwertig, die nach dem Modul $p-1$ kongruent sind, ebenso wie die Zahlen selbst als gleichwertig gelten, wenn sie nach dem Modul $p$ kongruent sind.

\Needspace{6\baselineskip}
\begin{center}\emph{Rechnen mit Indizes.}\end{center}

\Article{58}
Die Sätze über Indizes entsprechen genau denen über Logarithmen.

\emph{Der Index eines Produkts beliebig vieler Faktoren ist nach dem Modul $p-1$ der Summe der Indizes der einzelnen Faktoren kongruent.}

\emph{Der Index einer Potenz einer Zahl ist nach dem Modul $p-1$ dem Produkt aus dem Index der gegebenen Zahl und dem Exponenten der Potenz kongruent.}

Wegen ihrer Einfachheit lassen wir die Beweise weg.

Will man also eine Tabelle erstellen, aus der die Indizes aller Zahlen für verschiedene Moduln entnommen werden können, so darf man sowohl alle Zahlen über dem Modul als auch alle zusammengesetzten Zahlen weglassen. Ein Beispiel einer solchen Tabelle ist am Ende dieses Werkes als \emph{Tabelle} I beigefügt. In ihrer ersten senkrechten Spalte stehen die Primzahlen und Primzahlpotenzen von $3$ bis $97$, die als Moduln zu betrachten sind; daneben stehen die jeweils als Basis gewählten Zahlen. Es folgen die Indizes der aufeinanderfolgenden Primzahlen, die jeweils zu fünft durch einen kleinen
Zwischenraum\SourceLeaf{53} getrennt sind. Darüber stehen die Primzahlen in derselben Reihenfolge, sodass sich der Index einer gegebenen Primzahl nach einem gegebenen Modul leicht und zuverlässig finden lässt.

So ist \emph{zum Beispiel} für $p=67$ der Index von $60$, wenn $12$ als Basis gewählt wird,
\[\equiv2\operatorname{Ind.}2+\operatorname{Ind.}3+\operatorname{Ind.}5\ (\operatorname{mod.}66)\equiv58+9+39\equiv40\]

\Article{59}
\emph{Der Index eines beliebigen Wertes des Ausdrucks $\frac{a}{b}(\mathit{mod.}p.),$ (Artikel 31) ist nach dem Modul $p-1$ der Differenz der Indizes des Zählers $a$ und des Nenners $b$ kongruent, sofern die Zahlen $a,b$ nicht durch $p$ teilbar sind.}

Ein beliebiger Wert sei $c$. Dann ist $bc\equiv a\,(\operatorname{mod.}p)$, also
\[\operatorname{Ind.}b+\operatorname{Ind.}c\equiv\operatorname{Ind.}a\,(\operatorname{mod.}p-1)\]
und damit
\[\operatorname{Ind.}c\equiv\operatorname{Ind.}a-\operatorname{Ind.}b\]

Hat man eine Tabelle, die zu jeder Zahl für jeden Primzahlmodul den Index angibt, und eine weitere, die zu einem gegebenen Index die zugehörige Zahl liefert, so lassen sich alle Kongruenzen ersten Grades sehr leicht lösen. Denn alle lassen sich auf solche mit Primzahlmodul zurückführen (Artikel 30).

\emph{Zum Beispiel} ergebe sich aus der gegebenen Kongruenz
\[29x+7\equiv0\,(\operatorname{mod.}47)\quad\text{so ist}\quad x\equiv\frac{-7}{29}(\operatorname{mod.}47)\]
also
\[\operatorname{Ind.}x\equiv\operatorname{Ind.}-7-\operatorname{Ind.}29\equiv\operatorname{Ind.}40-\operatorname{Ind.}29\equiv15-43\equiv18(\operatorname{mod.}46)\]
Die Zahl mit dem Index $18$ ist $3$. Daher ist $x\equiv3\,(\operatorname{mod.}47)$. --- Die zweite Tabelle haben wir zwar nicht beigefügt; eine andere kann aber an ihre Stelle treten, wie wir in Abschnitt VI zeigen werden.

\Needspace{6\baselineskip}
\begin{center}\emph{Wurzeln der Kongruenz $x^n\equiv A$.}\end{center}

\Article{60}
Wie wir in Artikel 31 die Wurzeln der Kongruenzen ersten Grades bezeichnet haben, verwenden wir im Folgenden auch ein Zeichen für die Wurzeln reiner Kongruenzen höheren Grades. So wie ${}^{n}\!\surd A$ lediglich eine Wurzel der Gleichung $x^n=A$ bezeichnet, steht nach Hinzufügen des Moduls ${}^{n}\!\surd A\,(\operatorname{mod.}p)$ für jede beliebige Wurzel der Kongruenz $x^n\equiv A\,(\operatorname{mod.}p)$. Wir sagen, der Ausdruck ${}^{n}\!\surd A\,(\operatorname{mod.}p)$ habe so viele Werte, wie er nach $p$ inkongruente Werte hat. Denn alle nach $p$
\SourceLeaf{54}
kongruenten Werte sind als gleichwertig anzusehen (Artikel 26). Sind $A,B$ nach $p$ kongruent, so sind die Ausdrücke ${}^{n}\!\surd A,{}^{n}\!\surd B\,(\operatorname{mod.}p)$ offenbar gleichwertig.

Setzt man nun ${}^{n}\!\surd A\equiv x\,(\operatorname{mod.}p)$, so ist $n\operatorname{Ind.}x\equiv\operatorname{Ind.}A\,(\operatorname{mod.}p-1)$. Aus dieser Kongruenz erhält man nach den Regeln des vorigen Abschnitts die Werte von $\operatorname{Ind.}x$ und aus ihnen die entsprechenden Werte von $x$. Dabei hat $x$ offenbar ebenso viele Werte, wie die Kongruenz $n\operatorname{Ind.}x\equiv\operatorname{Ind.}A\,(\operatorname{mod.}p-1)$ Wurzeln hat. Somit hat ${}^{n}\!\surd A$ nur einen Wert, wenn $n$ zu $p-1$ teilerfremd ist. Haben dagegen die Zahlen $n,p-1$ den größten gemeinsamen Teiler $\delta$, so hat $\operatorname{Ind.}x$ genau $\delta$ nach $p-1$ inkongruente Werte und damit ${}^{n}\!\surd A$ ebenso viele nach $p$ inkongruente Werte, sofern $\operatorname{Ind.}A$ durch $\delta$ teilbar ist. Gilt diese Bedingung nicht, so hat ${}^{n}\!\surd A$ keinen wirklichen Wert.

\emph{Beispiel.} Gesucht sind die Werte des Ausdrucks ${}^{15}\!\surd11\,(\operatorname{mod.}19)$. Man muss also die Kongruenz $15\operatorname{Ind.}x\equiv\operatorname{Ind.}11\equiv6\,(\operatorname{mod.}18)$ lösen und findet die drei Werte $\operatorname{Ind.}x\equiv4,10,16\,(\operatorname{mod.}18)$. Ihnen entsprechen die Werte von $x$, nämlich $6,9,4$.
